% Source: https://artofproblemsolving.com/community/c3h2280235_problems__1000th_posts
% Stable identity: aops-topic-2280235-q05
Call the expression $E(a,b,c)$. Putting $a=-b$ gives $c^5-(-b)^5-b^5-c^5=0$, so, viewed as a polynomial in $a$, $E$ vanishes at $a=-b$ and by the factor theorem is divisible by $a+b$. By symmetry it is also divisible by $b+c$ and $c+a$, and these three factors are distinct linear polynomials, so
\[E=(a+b)(b+c)(c+a)\,Q(a,b,c),\]
where $Q$ is symmetric and homogeneous of degree $2$: $Q=k(a^2+b^2+c^2)+m(ab+bc+ca)$.
\begin{itemize}
\item $a=b=1$, $c=0$: $E=32-2=30$ and $(a+b)(b+c)(c+a)=2$, so $2k+m=15$.
\item $a=b=c=1$: $E=243-3=240$ and $(a+b)(b+c)(c+a)=8$, so $3k+3m=30$, i.e. $k+m=10$.
\end{itemize}
Hence $k=m=5$, and
\[(a+b+c)^5-a^5-b^5-c^5=\boxed{5(a+b)(b+c)(c+a)(a^2+b^2+c^2+ab+bc+ca)}.\]
(The quadratic factor does not factor further over the reals: it equals $\dfrac{1}{2}\left[(a+b)^2+(b+c)^2+(c+a)^2\right]$, which is positive unless $a=b=c=0$.)
